
\documentclass[11pt,a4paper]{article}
\usepackage[utf8]{inputenc}
\usepackage[T1]{fontenc}
\usepackage[spanish,es-noshorthands,es-tabla,es-nodecimaldot]{babel}
\usepackage{lmodern}


\usepackage[a4paper,top=2.4cm,bottom=2.4cm,left=2.3cm,right=2.3cm,
            headheight=24pt,footskip=1.1cm]{geometry}

\usepackage{amsmath,amssymb}
\usepackage{bm}
\usepackage{graphicx}
\usepackage{xcolor}
\usepackage{booktabs}
\usepackage{tabularx}
\usepackage{array}
\usepackage{longtable}
\usepackage{enumitem}
\usepackage{fancyhdr}
\usepackage{titlesec}
\usepackage{listings}
\usepackage[most]{tcolorbox}
\usepackage{hyperref}

\definecolor{navy}{HTML}{1E2761}
\definecolor{coral}{HTML}{F96167}
\definecolor{iceblue}{HTML}{CADCFC}
\definecolor{grisclaro}{HTML}{F4F6FB}

\hypersetup{
  colorlinks=true,
  linkcolor=navy,
  urlcolor=coral,
  citecolor=navy,
  pdftitle={Laboratorio 1 - Logica CMOS estatica},
  pdfauthor={MsC. Luz Adanaque Infante - UNMSM FIEE}
}

%------------------------------ Encabezados ----------------------------------
\pagestyle{fancy}
\fancyhf{}
\fancyhead[L]{\footnotesize\color{navy}\textsc{Microelectrónica} · Laboratorio N.\textsuperscript{o} 1}
\fancyhead[R]{\footnotesize\color{navy}Lógica CMOS estática}
\fancyfoot[L]{\footnotesize\color{navy}MsC. Luz Adanaqué Infante}
\fancyfoot[C]{\footnotesize\color{navy}UNMSM -- FIEE}
\fancyfoot[R]{\footnotesize\color{navy}\thepage}
\renewcommand{\headrulewidth}{0.8pt}
\renewcommand{\footrulewidth}{0.4pt}
\renewcommand{\headrule}{\hbox to\headwidth{\color{coral}\leaders\hrule height \headrulewidth\hfill}}

%------------------------------ Titulos --------------------------------------
\titleformat{\section}
  {\normalfont\Large\bfseries\color{navy}}{\thesection.}{0.6em}{}
  [\vspace{-0.6em}{\color{coral}\rule{\textwidth}{1.2pt}}]
\titleformat{\subsection}
  {\normalfont\large\bfseries\color{navy}}{\thesubsection}{0.6em}{}
\titlespacing*{\section}{0pt}{1.6em}{0.9em}

%------------------------- Estilo para codigo SPICE --------------------------
\lstdefinelanguage{SPICE}{
  alsoletter={.},
  morekeywords={.include,.param,.model,.tran,.dc,.ac,.op,.meas,.measure,.step,
                .ic,.subckt,.ends,.end,.options,.lib,.save,.four,
                PULSE,PWL,SIN,TRIG,TARG,VAL,RISE,FALL,TD,AVG,MAX,MIN,PARAM,
                FIND,WHEN,AT,DERIV,nmos,pmos,level},
  sensitive=false,
  morecomment=[f]{*},
  morecomment=[l]{;},
}
\lstdefinestyle{spice}{
  language=SPICE,
  basicstyle=\ttfamily\small,
  keywordstyle=\color{navy}\bfseries,
  commentstyle=\color{coral!85!black}\itshape,
  backgroundcolor=\color{grisclaro},
  frame=leftline,
  framerule=2pt,
  rulecolor=\color{iceblue},
  framesep=8pt,
  xleftmargin=10pt,
  breaklines=true,
  showstringspaces=false,
  columns=fullflexible,
  aboveskip=1em,
  belowskip=1em,
}
\lstset{style=spice}

%------------------------------- Cajas ---------------------------------------
\newtcolorbox{aviso}[1][]{
  enhanced, breakable, colback=coral!8, colframe=coral,
  boxrule=0pt, leftrule=3.5pt, arc=2pt, left=8pt, right=8pt, top=6pt, bottom=6pt,
  fonttitle=\bfseries\color{white}, #1}

\newtcolorbox{nota}[1][]{
  enhanced, breakable, colback=iceblue!35, colframe=navy,
  boxrule=0pt, leftrule=3.5pt, arc=2pt, left=8pt, right=8pt, top=6pt, bottom=6pt,
  #1}

\newtcolorbox{preguntas}[1][]{
  enhanced, breakable, colback=white, colframe=navy,
  boxrule=1pt, arc=3pt, left=8pt, right=8pt, top=6pt, bottom=6pt,
  title={\bfseries #1}, coltitle=white, colbacktitle=navy}

%------------------------- Contador de preguntas -----------------------------
\newcounter{preg}
\newcommand{\preg}[1]{\refstepcounter{preg}\item[\textbf{\color{coral}P\thepreg.}] #1}
\newenvironment{listapreg}{\begin{list}{}{%
  \setlength{\leftmargin}{2.4em}\setlength{\labelwidth}{2.0em}%
  \setlength{\itemsep}{0.5em}\setlength{\parsep}{0pt}}}{\end{list}}

%------------------------------ Utilidades -----------------------------------
\newcommand{\vacio}{\rule[-1.4ex]{0pt}{4.2ex}}
\renewcommand{\arraystretch}{1.35}
\setlength{\tabcolsep}{5pt}

%==============================================================================
\begin{document}
%==============================================================================

%------------------------------- Portada -------------------------------------
\begin{titlepage}
\centering
\vspace*{1.2cm}

{\color{navy}\rule{\textwidth}{2pt}}\\[0.5em]
{\Large\bfseries\color{navy} UNIVERSIDAD NACIONAL MAYOR DE SAN MARCOS}\\[0.35em]
{\normalsize\color{navy} Universidad del Perú, Decana de América}\\[0.6em]
{\color{coral}\rule{\textwidth}{2pt}}\\[1.4cm]


{\large\color{navy} Facultad de Ingeniería Electrónica y Eléctrica}\\[0.3em]
{\large\color{navy} Escuela Profesional de Ingeniería Electrónica}\\[2.2cm]

{\color{coral}\Large\bfseries GUÍA DE LABORATORIO N.\textsuperscript{o} 1}\\[0.9cm]
{\Huge\bfseries\color{navy} Lógica CMOS Estática}\\[0.7cm]
{\large\color{navy}\itshape Caracterización, dimensionamiento y potencia\\ con modelos predictivos PTM en LTspice}\\[2.4cm]

\begin{tabular}{@{}r@{\hspace{1em}}l@{}}
\bfseries\color{navy} Curso: & Microelectrónica \\

\bfseries\color{navy} Unidad: & Modelamiento de circuitos \\
\bfseries\color{navy} Docente: & MsC. Luz Adanaqué Infante \\
\bfseries\color{navy} Duración: & 4 horas de laboratorio\\
\bfseries\color{navy} Modalidad: & Parejas \\
\end{tabular}

\vfill
{\color{iceblue}\rule{\textwidth}{6pt}}\\[0.4em]
{\footnotesize\color{navy} Lima, Perú}
\end{titlepage}

%--------------------------------- Indice ------------------------------------
{\color{navy}\tableofcontents}
\newpage

%==============================================================================
\section{Objetivos}
%==============================================================================

Al terminar la sesión el estudiante será capaz de:

\begin{enumerate}[leftmargin=1.6em,itemsep=0.35em]
  \item Configurar un entorno de simulación con tarjetas de modelo predictivas (PTM) para distintos nodos tecnológicos.
  \item Extraer la curva de transferencia (VTC) de un inversor CMOS estático y determinar $V_M$, $V_{IL}$, $V_{IH}$ y los márgenes de ruido.

  \item Medir retardos de propagación, tiempos de transición y su dependencia con la carga y el nodo tecnológico.
  \item Sintetizar compuertas complejas (NAND, NOR, AOI) mediante redes pull-up / pull-down duales y verificarlas funcionalmente.
  \item Cuantificar las tres componentes de potencia (dinámica, cortocircuito y fuga) y contrastar procesos HP frente a LP.
\end{enumerate}

%==============================================================================
\section{Fundamento resumido}
%==============================================================================

La lógica CMOS estática implementa una función $Y = \overline{F(A,B,\dots)}$ mediante dos redes complementarias: una red \emph{pull-down} (PDN) de NMOS que conduce cuando $F=1$, y una red \emph{pull-up} (PUN) de PMOS, grafo dual de la anterior. Como en régimen estático solo una de las dos redes conduce, en el modelo ideal no existe camino DC entre $V_{DD}$ y tierra: la potencia estática nace únicamente de las corrientes de fuga.

Tres relaciones gobiernan el trabajo de hoy.

\paragraph{Punto de conmutación.} Igualando las corrientes de ambos transistores en $V_{in}=V_{out}=V_M$, con los dispositivos en saturación de velocidad:

\begin{equation}
V_M = \frac{V_{Tn} + r\,(V_{DD}+V_{Tp})}{1+r},
\qquad
r = \frac{k_p\,V_{DSATp}}{k_n\,V_{DSATn}} \simeq \frac{v_{sat,p}\,W_p}{v_{sat,n}\,W_n}
\label{eq:vm}
\end{equation}

\noindent donde $V_{Tp}<0$. En tecnologías de canal largo la razón $r$ se aproxima por $\mu_p W_p / \mu_n W_n$; por debajo de los 65~nm la saturación de velocidad domina y la razón de movilidades deja de ser el criterio correcto. Este contraste es el eje de la Actividad~1.

\paragraph{Retardo de propagación.}
\begin{equation}
t_p = \tfrac{1}{2}\left(t_{pHL}+t_{pLH}\right),
\qquad
t_{pHL} \approx 0.69\,R_{eq,n}\,C_L
\label{eq:tp}
\end{equation}

\paragraph{Potencia.}
\begin{equation}
P_{tot} = \underbrace{\alpha\,C_L V_{DD}^{2} f}_{\text{dinámica}}
        + \underbrace{I_{cc}\,V_{DD}}_{\text{cortocircuito}}
        + \underbrace{I_{leak}\,V_{DD}}_{\text{fuga}}
\label{eq:pot}
\end{equation}


\section{Preparación del entorno}


\subsection{Modelos disponibles}

Todos los modelos provienen del \emph{Predictive Technology Model} (Arizona State University). Los valores de la Tabla~\ref{tab:ptm} ya están extraídos de los archivos publicados en el aula virtual y sirven de referencia rápida.

\begin{table}[htbp]
\centering
\small
\caption{Tarjetas de modelo PTM disponibles en el aula virtual.}
\label{tab:ptm}
\begin{tabular}{@{}llccccl@{}}
\toprule
\textbf{Archivo} & \textbf{Nodo} & $\bm{V_{DD}}$ & $\bm{V_{TH0,n}}$ & $\bm{V_{TH0,p}}$ & $\bm{t_{oxe,n}}$ & \textbf{Modelo} \\
 & & (V) & (V) & (V) & (nm) & \\
\midrule
\texttt{180nm\_bulk.txt} & 180 nm    & 1.8\textsuperscript{*} & 0.3999  & $-0.42$    & 4.00 & BSIM3 (lvl 49) \\
\texttt{130nm\_bulk.pm}  & 130 nm    & 1.3\textsuperscript{*} & 0.3782  & $-0.321$   & 2.25 & BSIM4 (lvl 54) \\
\texttt{90nm\_bulk.pm}   & 90 nm     & 1.2\textsuperscript{*} & 0.397   & $-0.339$   & 2.05 & BSIM4 \\
\texttt{65nm\_bulk.pm}   & 65 nm     & 1.1\textsuperscript{*} & 0.423   & $-0.365$   & 1.85 & BSIM4 \\
\texttt{45nm\_HP.pm}     & 45 nm HP  & 1.0 & 0.46893 & $-0.49158$ & 1.25 & BSIM4 \\
\texttt{45nm\_LP.pm}     & 45 nm LP  & 1.1 & 0.62261 & $-0.587$   & 1.80 & BSIM4 \\
\texttt{32nm\_HP.pm}     & 32 nm HP  & 0.9 & 0.49396 & $-0.49155$ & 1.15 & BSIM4 \\
\texttt{32nm\_LP.pm}     & 32 nm LP  & 1.0 & 0.63    & $-0.5808$  & 1.60 & BSIM4 \\
\texttt{22nm\_HP.pm}     & 22 nm HP  & 0.8 & 0.50308 & $-0.4606$  & 1.05 & BSIM4 \\
\texttt{22nm\_LP.pm}     & 22 nm LP  & 0.95& 0.68858 & $-0.63745$ & 1.40 & BSIM4 \\
\bottomrule
\end{tabular}

\vspace{0.4em}
\parbox{0.92\textwidth}{\footnotesize \textsuperscript{*}Los cuatro procesos \emph{bulk} no declaran $V_{DD}$ dentro del archivo; use el valor referencial de la tabla y declárelo explícitamente en su netlist.}
\end{table}



\subsection{Inclusión de la tarjeta en LTspice}

\begin{enumerate}[leftmargin=1.6em,itemsep=0.35em]
  \item Copie el archivo \texttt{.pm} a la \textbf{misma carpeta} que su esquemático o netlist.
  \item Agregue como directiva SPICE (tecla \texttt{S} en el esquemático):
        \lstinline|.include 45nm\_HP.pm|
  \item En el esquemático use los símbolos de \textbf{cuatro terminales} \texttt{nmos4} y \texttt{pmos4} (D--G--S--B). Los de tres terminales asumen sustrato conectado a fuente y le impedirán estudiar el efecto de cuerpo en la Actividad~3.
  \item Haga clic derecho sobre cada transistor y escriba el nombre del modelo: \texttt{nmos} o \texttt{pmos} (así se llaman literalmente dentro de las tarjetas), además de \texttt{L} y \texttt{W}.
\end{enumerate}

\paragraph{Longitud de canal.} Use como longitud dibujada el valor nominal del nodo (\texttt{L=45n} para la tarjeta de 45~nm). La tarjeta ya incorpora \texttt{xl} y \texttt{lint} para obtener la longitud efectiva; \textbf{no} corrija usted esos valores.

\paragraph{Anchos.} Trabaje siempre en múltiplos del ancho mínimo. En esta guía se adopta $W_n = 4L$ como transistor unitario.

\subsection{Colisión de nombres al comparar tecnologías}
\label{sec:colision}

Todas las tarjetas definen sus modelos con los mismos nombres, \texttt{nmos} y \texttt{pmos}. Si incluye dos archivos en la misma simulación, el segundo sobrescribe al primero \textbf{sin emitir error}, y usted comparará dos veces la misma tecnología sin darse cuenta. Es el error más frecuente de este laboratorio.

La solución es renombrar los modelos antes de comparar. Edite una copia del archivo y cambie la primera línea de cada modelo:

\begin{lstlisting}
.model  nmos  nmos  level = 54     ->   .model nmos_45hp nmos level = 54
.model  pmos  pmos  level = 54     ->   .model pmos_45hp pmos level = 54
\end{lstlisting}

\noindent o, desde una terminal:

\begin{lstlisting}[language={}]
sed -E 's/^\.model[[:space:]]+nmos/.model nmos_45hp/;
        s/^\.model[[:space:]]+pmos/.model pmos_45hp/' 45nm_HP.pm > 45nm_HP_ren.lib
\end{lstlisting}

\subsection{Nota de compatibilidad}

Las tarjetas declaran \texttt{level = 54} (BSIM4). Si su versión de LTspice reporta un nivel desconocido, reemplace por \texttt{level = 14}, que es el identificador interno de BSIM4 en LTspice. Es normal que el registro de errores liste parámetros ignorados; mientras la simulación converja, esos avisos no invalidan los resultados.

%==============================================================================
\section{Actividad 1 --- Inversor CMOS: VTC y márgenes de ruido}
%==============================================================================

\subsection*{Procedimiento}
\begin{enumerate}[leftmargin=1.6em,itemsep=0.35em]
  \item Arme el circuito del Anexo~\ref{sec:a1} con la tarjeta de 45~nm HP, $V_{DD}=1.0$~V, $L=45$~nm, $W_n=180$~nm.
  \item Ejecute el barrido DC con \lstinline|.step param beta 1 3 0.5| y grafique la familia de curvas $V_{out}(V_{in})$.
  \item Registre $V_M$ para cada valor de $\beta$ mediante la directiva \lstinline|.meas|. Los resultados se leen en el archivo de registro (\texttt{Ctrl+L}).
  \item Para el caso de $\beta$ que hace $V_M \approx V_{DD}/2$: grafique la derivada \lstinline|d(V(out))| en el visor de formas de onda y ubique con cursores los dos puntos donde la pendiente vale $-1$. Esos son $V_{IL}$ y $V_{IH}$.
  \item Calcule $NM_H = V_{OH}-V_{IH}$ y $NM_L = V_{IL}-V_{OL}$.
\end{enumerate}

\begin{table}[htbp]
\centering\small
\caption{Resultados 1.1 --- Barrido de la razón $\beta$.}
\begin{tabularx}{\textwidth}{@{}c*{6}{>{\centering\arraybackslash}X}@{}}
\toprule
$\bm{\beta}$ & $\bm{V_M}$ (V) & $\bm{V_{IL}}$ (V) & $\bm{V_{IH}}$ (V) & $\bm{NM_L}$ (V) & $\bm{NM_H}$ (V) & \textbf{Gan. máx.} \\
\midrule
1.0 & \vacio & & & & & \\
1.5 & \vacio & & & & & \\
2.0 & \vacio & & & & & \\
2.5 & \vacio & & & & & \\
3.0 & \vacio & & & & & \\
\bottomrule
\end{tabularx}
\end{table}

\begin{table}[htbp]
\centering\small
\caption{Resultados 1.2 --- Parámetros extraídos de \texttt{45nm\_HP.pm}. Abra el archivo con un editor de texto y complete.}
\begin{tabularx}{0.85\textwidth}{@{}l*{3}{>{\centering\arraybackslash}X}@{}}
\toprule
\textbf{Parámetro} & \textbf{NMOS} & \textbf{PMOS} & \textbf{Razón n/p} \\
\midrule
\texttt{u0} (movilidad) & \vacio & & \\
\texttt{vsat}           & \vacio & & \\
\texttt{vth0}           & \vacio & & \\
\bottomrule
\end{tabularx}
\end{table}

\begin{preguntas}[Preguntas de la Actividad 1]
\begin{listapreg}
\preg{Con la razón de \texttt{u0} obtenida, la regla clásica predice cierto $\beta$ para centrar la VTC. ¿Coincide con el $\beta$ que usted midió? Explique la discrepancia en términos de saturación de velocidad, apoyándose en la ecuación~\eqref{eq:vm}.}
\preg{¿Por qué $V_{OH}$ y $V_{OL}$ son exactamente $V_{DD}$ y 0~V en lógica CMOS estática, y no lo son en lógica pseudo-NMOS?}
\preg{Un $\beta$ mayor mejora un margen de ruido y degrada el otro. ¿Cuál y por qué? ¿En qué aplicación preferiría un inversor deliberadamente descentrado?}
\preg{Repita el barrido con \texttt{45nm\_LP.pm} a $V_{DD}=1.1$~V. ¿Cómo cambia la ganancia en la región de transición y qué parámetro de la tarjeta lo explica?}
\end{listapreg}
\end{preguntas}

%==============================================================================
\section{Actividad 2 --- Respuesta transitoria: retardo y transición}
%==============================================================================

\subsection*{Procedimiento}
\begin{enumerate}[leftmargin=1.6em,itemsep=0.35em]
  \item Use el netlist del Anexo~\ref{sec:a2}, con el $\beta$ óptimo de la Actividad~1.
  \item Mida $t_{pHL}$, $t_{pLH}$, $t_r$ y $t_f$ con las directivas \lstinline|.meas| provistas.
  \item Barra la carga con \lstinline|.step param Cload 1f 8f 1f| y grafique $t_p$ contra $C_L$. Ajuste una recta y determine su pendiente y su ordenada al origen.
  \item Repita la medición para 90~nm, 45~nm HP y 22~nm HP, escalando $V_{DD}$ y $L$ según la Tabla~\ref{tab:ptm} y manteniendo $W_n=4L$ y $C_L=2$~fF.
\end{enumerate}

\begin{table}[htbp]
\centering\small
\caption{Resultados 2.1 --- Retardo frente al nodo tecnológico.}
\begin{tabularx}{\textwidth}{@{}lc*{5}{>{\centering\arraybackslash}X}@{}}
\toprule
\textbf{Nodo} & $\bm{V_{DD}}$ (V) & $\bm{t_{pHL}}$ (ps) & $\bm{t_{pLH}}$ (ps) & $\bm{t_p}$ (ps) & $\bm{t_r}$ (ps) & $\bm{t_f}$ (ps) \\
\midrule
90 nm     & 1.2 & \vacio & & & & \\
45 nm HP  & 1.0 & \vacio & & & & \\
22 nm HP  & 0.8 & \vacio & & & & \\
\bottomrule
\end{tabularx}
\end{table}

\begin{preguntas}[Preguntas de la Actividad 2]
\begin{listapreg}
\preg{De la recta $t_p$ contra $C_L$: ¿qué representa físicamente la pendiente y qué representa la ordenada al origen?}
\preg{Con $\beta$ correctamente ajustado, $t_{pHL}$ y $t_{pLH}$ deberían ser casi iguales. Si en su simulación difieren en más de 10~\%, ¿qué corrección aplicaría y por qué el ajuste no resulta idéntico al que centró la VTC?}
\preg{Entre 90~nm y 22~nm el retardo mejora aunque $V_{DD}$ se reduce en un tercio. Identifique los dos mecanismos del escalamiento que compensan la caída de tensión.}
\preg{Aumente el tiempo de subida del estímulo de 5~ps a 100~ps. ¿Qué componente de la ecuación~\eqref{eq:pot} crece y por qué?}
\end{listapreg}
\end{preguntas}

\newpage

%==============================================================================
\section{Actividad 3 --- NAND2 y NOR2}
%==============================================================================

\subsection*{Procedimiento}
\begin{enumerate}[leftmargin=1.6em,itemsep=0.35em]
  \item Implemente la NAND2 (Anexo~\ref{sec:a3}) y la NOR2 con el criterio de \textbf{resistencia equivalente igual a la del inversor de referencia}: en la rama serie los transistores duplican su ancho; en la rama paralela conservan el ancho del inversor.
  \item Verifique la tabla de verdad de cada compuerta.
  \item Mida $t_{pHL}$ de la NAND2 en dos escenarios: (a) conmuta la entrada \textbf{interna} (la conectada al transistor cercano a tierra) con la otra fija en \texttt{1}; (b) conmuta la entrada \textbf{externa} con la otra fija en \texttt{1}.
  \item Mida $t_{pLH}$ de la NOR2 con las dos entradas conmutando simultáneamente y con solo una conmutando.
  \item Compare NAND2 y NOR2 con idéntica carga: retardo promedio y ancho total de transistores (proxy de área).
\end{enumerate}

\begin{table}[htbp]
\centering\small
\caption{Resultados 3.1 --- Comparación de compuertas.}
\begin{tabularx}{\textwidth}{@{}l*{6}{>{\centering\arraybackslash}X}@{}}
\toprule
\textbf{Compuerta} & $\bm{W_n}$/trans. & $\bm{W_p}$/trans. & $\bm{\sum W}$ & $\bm{t_{pHL}}$ (ps) & $\bm{t_{pLH}}$ (ps) & $\bm{t_p}$ (ps) \\
\midrule
Inversor (ref.) & \vacio & & & & & \\
NAND2           & \vacio & & & & & \\
NOR2            & \vacio & & & & & \\
\bottomrule
\end{tabularx}
\end{table}

\begin{preguntas}[Preguntas de la Actividad 3]
\begin{listapreg}
\preg{Los dos escenarios del paso 3 dan retardos distintos. Explique el papel de la capacitancia del nodo intermedio y del efecto de cuerpo del transistor superior.}
\preg{El esfuerzo lógico de la NAND2 es $4/3$ y el de la NOR2 es $5/3$. Verifique cualitativamente esa jerarquía con sus mediciones y explique por qué la industria construye bibliotecas dominadas por NAND.}
\preg{¿Qué le ocurriría al retardo de una NOR de cuatro entradas dimensionada con este mismo criterio? Estime el $W_p$ requerido y comente su viabilidad.}
\preg{Si el sustrato del NMOS superior de la NAND2 se conectara a su fuente en lugar de a tierra, ¿qué cambiaría en el retardo? ¿Es esto físicamente realizable en un proceso \emph{bulk}?}
\end{listapreg}
\end{preguntas}

%==============================================================================
\section{Actividad 4 --- Síntesis de una compuerta compleja}
%==============================================================================

Diseñe, dimensione y verifique la compuerta que implementa
\begin{equation}
Y = \overline{(A \cdot B) + C}
\end{equation}

\subsection*{Procedimiento}
\begin{enumerate}[leftmargin=1.6em,itemsep=0.35em]
  \item Dibuje la PDN a partir de la expresión (serie $\leftrightarrow$ AND, paralelo $\leftrightarrow$ OR) y obtenga la PUN como grafo dual.
  \item Dimensione cada rama para igualar la resistencia equivalente del inversor de referencia. Justifique cada ancho por escrito.
  \item Verifique las ocho combinaciones de entrada usando el estímulo binario del Anexo~\ref{sec:a4} ($C$ con período $T$, $B$ con $2T$, $A$ con $4T$).
  \item Identifique en la simulación la combinación de entrada que produce el peor caso de $t_{pHL}$ y la que produce el peor caso de $t_{pLH}$.
\end{enumerate}

\begin{preguntas}[Preguntas de la Actividad 4]
\begin{listapreg}
\preg{¿Cuántos transistores requiere su implementación? Compárela con la realización de la misma función mediante NAND2 más inversores en cascada, en número de transistores y en número de etapas.}
\preg{La PUN de su diseño tiene una rama serie de dos PMOS. ¿Por qué esta función es más ``amigable'' para CMOS estático que su complemento $Y=\overline{(A+B)\cdot C}$?}
\preg{Explique por qué la lógica CMOS estática siempre produce una función \textbf{inversora} y qué implica esto al mapear una expresión booleana arbitraria.}
\end{listapreg}
\end{preguntas}

%==============================================================================
\section{Actividad 5 --- Potencia: dinámica, cortocircuito y fuga}
%==============================================================================

\subsection*{Procedimiento}
\begin{enumerate}[leftmargin=1.6em,itemsep=0.35em]
  \item \textbf{Fuga.} Con el netlist~\ref{sec:a1} modificado para punto de operación: fije \texttt{Vin~=~0} y luego \texttt{Vin~=~VDD}, y registre $|I(V1)|$ en cada caso. Repita con \texttt{45nm\_HP.pm} y \texttt{45nm\_LP.pm}.
  \item \textbf{Dinámica.} En el netlist~\ref{sec:a2}, use \lstinline|.meas tran Iavg AVG I(V1)| y calcule $P=|I_{avg}|\,V_{DD}$. Verifique contra $\alpha C_L V_{DD}^2 f$ y explique cualquier exceso.
  \item \textbf{Oscilador de anillo.} Arme el anillo de cinco etapas del Anexo~\ref{sec:a5}. Mida el período de oscilación, deduzca $t_{pd}=T_{osc}/(2N)$ y compárelo con el $t_p$ de la Actividad~2.
  \item Repita el anillo con la tarjeta LP del mismo nodo y complete la tabla.
\end{enumerate}

\begin{table}[htbp]
\centering\small
\caption{Resultados 5.1 --- Compromiso velocidad/fuga entre procesos HP y LP.}
\begin{tabularx}{\textwidth}{@{}lc*{4}{>{\centering\arraybackslash}X}@{}}
\toprule
\textbf{Proceso} & $\bm{V_{DD}}$ (V) & $\bm{I_{leak}}$ (A) & $\bm{f_{osc}}$ (GHz) & $\bm{t_{pd}}$ (ps) & $\bm{P_{prom}}$ (\textmu W) \\
\midrule
45 nm HP & 1.0 & \vacio & & & \\
45 nm LP & 1.1 & \vacio & & & \\
\bottomrule
\end{tabularx}
\end{table}

\begin{preguntas}[Preguntas de la Actividad 5]
\begin{listapreg}
\preg{Calcule las razones $I_{leak,HP}/I_{leak,LP}$ y $f_{osc,HP}/f_{osc,LP}$. Enuncie el compromiso de diseño en una frase y cite el parámetro de la tarjeta que lo origina.}
\preg{La corriente de fuga que midió no proviene de un solo mecanismo. Nombre al menos tres e indique con qué parámetros del modelo BSIM4 se controlan en las tarjetas PTM (revise \texttt{voff}, \texttt{nfactor}, \texttt{eta0}, \texttt{agidl}, \texttt{aigc}).}
\preg{¿Por qué el $t_{pd}$ del oscilador de anillo suele ser menor que el $t_p$ medido con un pulso ideal y carga fija?}
\preg{Para un procesador móvil que pasa 95~\% del tiempo en reposo, ¿qué proceso elegiría y qué técnica adicional aplicaría al 5~\% activo?}
\preg{La potencia de cortocircuito no aparece explícitamente en sus mediciones porque está incluida en $I_{avg}$. Proponga un procedimiento de simulación que permita separarla de la componente de carga y descarga.}
\end{listapreg}
\end{preguntas}

%==============================================================================
\section{Entregables}
%==============================================================================

Informe en PDF, máximo 12 páginas, con:

\begin{enumerate}[leftmargin=1.6em,itemsep=0.3em]
  \item Portada, integrantes y declaración de autoría.
  \item Pre-lab resuelto: deducción de la ecuación~\eqref{eq:vm} y esquemas.
  \item Tablas 1.1, 1.2, 2.1, 3.1 y 5.1 completas.
  \item Gráficas: familia de VTC, $t_p$ contra $C_L$, transitorio de verificación funcional de la Actividad~4 y oscilación del anillo.
  \item Respuestas a las \thepreg{} preguntas, cada una en un máximo de cinco líneas.
  \item Anexo con los netlists finales en texto plano, no capturas de pantalla.
  \item Archivos \texttt{.asc} / \texttt{.cir} en un ZIP separado.
\end{enumerate}

\begin{nota}
\textbf{Entrega: 11 de setiembre a la medianoche} \\
\textbf{Consultas:} \href{mailto:Ladanaquei@unmsm.edu.pe}{Ladanaquei@unmsm.edu.pe}
\end{nota}

%==============================================================================
\section{Rúbrica de evaluación}
%==============================================================================

\begin{table}[htbp]
\centering\small
\caption{Rúbrica sobre 20 puntos.}
\begin{tabularx}{\textwidth}{@{}Xc@{}}
\toprule
\textbf{Criterio} & \textbf{Puntos} \\
\midrule
Pre-lab completo y correcto & 2 \\
Actividad 1: extracción de VTC y márgenes de ruido & 3 \\
Actividad 2: caracterización de retardo y escalamiento & 3 \\
Actividad 3: dimensionamiento y análisis del patrón de entrada & 3 \\
Actividad 4: síntesis y verificación funcional & 3 \\
Actividad 5: potencia y comparación HP/LP & 3 \\
Calidad del informe: gráficas rotuladas, unidades, discusión & 2 \\
Reproducibilidad: netlists ejecutables sin corrección & 1 \\
\midrule
\textbf{Total} & \textbf{20} \\
\bottomrule
\end{tabularx}
\end{table}

\noindent\textbf{Descuentos.} $-2$ puntos si las gráficas carecen de ejes rotulados con unidades; $-3$ puntos si se comparan tecnologías sin renombrar los modelos (sección~\ref{sec:colision}), pues los resultados serían inválidos.

%==============================================================================
\appendix
\section{Netlists de referencia}
%==============================================================================

Guarde cada bloque como archivo \texttt{.cir} y ejecútelo con \texttt{File $\rightarrow$ Open} en LTspice, o replique el circuito en esquemático.

\subsection{Curva de transferencia del inversor}
\label{sec:a1}

\begin{lstlisting}
* === A1: VTC del inversor CMOS - PTM 45nm HP ===
.include 45nm_HP.pm

.param VDDv = 1.0
.param Lmin = 45n
.param Wn   = 4*Lmin
.param beta = 2

V1  vdd 0 {VDDv}
VIN in  0 0

MN out in 0   0   nmos L={Lmin} W={Wn}
MP out in vdd vdd pmos L={Lmin} W={beta*Wn}

.dc VIN 0 {VDDv} 2m
.step param beta 1 3 0.5
.meas dc VM FIND V(in) WHEN V(out)=V(in)
.end
\end{lstlisting}

Para el punto de fuga de la Actividad~5, reemplace las dos últimas directivas de análisis por:

\begin{lstlisting}
.dc VIN 0 {VDDv} {VDDv}
.meas dc Ileak_lo FIND I(V1) AT 0
.meas dc Ileak_hi FIND I(V1) AT {VDDv}
\end{lstlisting}

\subsection{Respuesta transitoria}
\label{sec:a2}

\begin{lstlisting}
* === A2: Retardo de propagacion del inversor ===
.include 45nm_HP.pm

.param VDDv  = 1.0
.param Lmin  = 45n
.param Wn    = 4*Lmin
.param beta  = 2
.param Cload = 2f

V1  vdd 0 {VDDv}
VIN in  0 PULSE(0 {VDDv} 200p 5p 5p 500p 1n)

MN out in 0   0   nmos L={Lmin} W={Wn}
MP out in vdd vdd pmos L={Lmin} W={beta*Wn}
C1  out 0 {Cload}

.tran 0 2n 0 0.5p

.meas tran tpHL TRIG V(in)  VAL={VDDv/2} TD=100p RISE=1 TARG V(out) VAL={VDDv/2} FALL=1
.meas tran tpLH TRIG V(in)  VAL={VDDv/2} TD=100p FALL=1 TARG V(out) VAL={VDDv/2} RISE=1
.meas tran tf   TRIG V(out) VAL={0.9*VDDv} FALL=1 TARG V(out) VAL={0.1*VDDv} FALL=1
.meas tran tr   TRIG V(out) VAL={0.1*VDDv} RISE=1 TARG V(out) VAL={0.9*VDDv} RISE=1
.meas tran tp   PARAM (tpHL+tpLH)/2
.meas tran Iavg AVG I(V1)

.step param Cload 1f 8f 1f
.end
\end{lstlisting}

\subsection{NAND2 estática}
\label{sec:a3}

\begin{lstlisting}
* === A3: NAND2 CMOS estatica ===
.include 45nm_HP.pm

.param VDDv = 1.0
.param Lmin = 45n
.param Wn   = 4*Lmin
.param beta = 2

V1 vdd 0 {VDDv}
VA a 0 PULSE(0 {VDDv} 200p 10p 10p 1n 2n)
VB b 0 {VDDv}

* --- PDN: NMOS en serie, ancho duplicado ---
MN1 out a nx  0   nmos L={Lmin} W={2*Wn}    ; entrada externa
MN2 nx  b 0   0   nmos L={Lmin} W={2*Wn}    ; entrada interna
* --- PUN: PMOS en paralelo, ancho del inversor ---
MP1 out a vdd vdd pmos L={Lmin} W={beta*Wn}
MP2 out b vdd vdd pmos L={Lmin} W={beta*Wn}

C1 out 0 2f

.tran 0 4n 0 0.5p
.meas tran tpHL TRIG V(a) VAL={VDDv/2} RISE=1 TARG V(out) VAL={VDDv/2} FALL=1
.end
\end{lstlisting}

Para el escenario de entrada interna, intercambie los estímulos (\lstinline|VA a 0 {VDDv}| y aplique el pulso a \texttt{b}), ajustando el disparo de \lstinline|.meas|. Para la NOR2, invierta la topología: NMOS en paralelo con ancho \lstinline|{Wn}| y PMOS en serie con ancho \lstinline|{2*beta*Wn}|.

\subsection{Compuerta compleja: verificación funcional}
\label{sec:a4}

Estímulo binario para recorrer las ocho combinaciones, con $T=1$~ns:

\begin{lstlisting}
VA a 0 PULSE(0 {VDDv} 0 10p 10p 4n 8n)
VB b 0 PULSE(0 {VDDv} 0 10p 10p 2n 4n)
VC c 0 PULSE(0 {VDDv} 0 10p 10p 1n 2n)
.tran 0 8n 0 1p
\end{lstlisting}

Complete usted mismo la PDN y la PUN según su diseño de la Actividad~4, y agregue \lstinline|C1 out 0 2f|.

\subsection{Oscilador de anillo de cinco etapas}
\label{sec:a5}

\begin{lstlisting}
* === A5: Oscilador de anillo, 5 etapas ===
.include 45nm_HP.pm

.param VDDv = 1.0
.param Lmin = 45n
.param Wn   = 4*Lmin
.param beta = 2
.param N    = 5

.subckt INV in out vdd
MN out in 0   0   nmos L={Lmin} W={Wn}
MP out in vdd vdd pmos L={Lmin} W={beta*Wn}
.ends

V1 vdd 0 {VDDv}
X1 n1 n2 vdd INV
X2 n2 n3 vdd INV
X3 n3 n4 vdd INV
X4 n4 n5 vdd INV
X5 n5 n1 vdd INV

.ic V(n1)=0
.tran 0 20n 5n 0.5p

.meas tran Tosc TRIG V(n1) VAL={VDDv/2} RISE=2 TARG V(n1) VAL={VDDv/2} RISE=3
.meas tran fosc PARAM 1/Tosc
.meas tran tpd  PARAM Tosc/(2*N)
.meas tran Iavg AVG I(V1)
.end
\end{lstlisting}

%==============================================================================
\section{Problemas frecuentes}
%==============================================================================

\begin{table}[htbp]
\centering\small
\caption{Diagnóstico rápido de fallas en LTspice.}
\begin{tabularx}{\textwidth}{@{}p{4.3cm}XX@{}}
\toprule
\textbf{Síntoma} & \textbf{Causa probable} & \textbf{Solución} \\
\midrule
\texttt{Unknown model}, o el transistor aparece sin parámetros
& El archivo \texttt{.pm} no está en la carpeta del esquemático, o el nombre del modelo está mal escrito
& Verifique la ruta del \texttt{.include} y que el nombre sea exactamente \texttt{nmos} / \texttt{pmos} \\
Dos tecnologías dan resultados idénticos
& Colisión de nombres de modelo
& Renombre los modelos (sección~\ref{sec:colision}) \\
El oscilador de anillo no arranca
& Todos los nodos quedaron en el punto metaestable
& Agregue \texttt{.ic V(n1)=0} \\
\texttt{Timestep too small}
& Paso de simulación grueso o carga nula
& Reduzca el paso máximo del \texttt{.tran} y verifique que exista \texttt{C1} en la salida \\
Corriente de fuga con signo negativo
& Convención de LTspice para fuentes de tensión
& Use el valor absoluto: la corriente entra por el terminal positivo \\
Retardos absurdos (ns en lugar de ps)
& Se omitieron los sufijos de unidad
& Revise que \texttt{W} y \texttt{L} lleven sufijo \texttt{n} \\
\bottomrule
\end{tabularx}
\end{table}

\end{document}